\chapter{Novo capítulo}

\noindent{}Em \emph{Borderland/\,La Frontera}, a poeta e ensaísta Gloria Anzaldúa diz
que os grupos fronteiriços vivem no entre, em \emph{nepantla}:
``nepantla é estar em meio a uma transformação {[}\ldots{}{]} nepantla
é um jeito de ler o mundo''\footnote{\textsc{anzaldúa}, \emph{op.\,cit.}, p.\,237. Tradução da autora.}. Viver no entre --- no limiar --- é uma condição das
comunidades e da identidade judaica. Estar em fronteira é estar
diretamente ligado à alteridade, em contato e conexão com os demais
grupos sociais que se estabelecem nestes territórios.

Em \emph{Limiar, aura e rememoração}, a pesquisadora Jeanne Marie
Gagnebin argumenta que estes limiares são zonas de intermédio e rituais
de passagem que permitem suspender o espaço, e assim fugir de definições
conclusivas e autocircundantes, tais como: público/\,privado, vida/\,morte,
sagrado/\,profano. A experiência do limiar pode ser lida enquanto uma
soleira ou o umbral de uma porta (da mesma maneira que é o \emph{eruv}),
uma zona de transição e de mudança.

Atravessando fronteiras, os povos em diáspora contestam e questionam as
conotações. E, dessa forma, o viver no limiar permite, justamente, a
abertura de poros para o movimento de não apenas ocupar, mas sim de
habitar o entre. Este espaço entre, finalmente, possibilita
``reconquistar para o pensamento os territórios do indeterminado e
do intermediário, da suspensão e da hesitação''\footnote{GAGNEBIN,
  \emph{Limiar, aura e rememoração}, p.\,40.}.

Ao habitar a diáspora, a identidade judaica e as agregações
comunitárias, por meio de suas tradições, cultos e ritos, bem como de
seus idiomas e de suas línguas, se apropriam dos espaços de trânsito ---
estes movimentos de dispersão cíclicos e constantes --- e os tornam uma
possibilidade. Em outras palavras, as comunidades judaicas (e seus
indivíduos) transformam essas condições de existência --- estar entre
espaços --- em um local (e um tempo) de pertencimento.

Os movimentos apropriativos da condição diaspórica assemelham-se a
ecótonos. Ecótonos são zonas de transição ambiental, regiões resultantes
do contato entre dois ou mais biomas fronteiriços, em que se encontram
as mais diversas comunidades ecológicas. São territórios ricos em
espécies, pois compartilham espaços e permitem que sejam influenciados
por outros ambientes. É exatamente desta maneira que devem ser lidas as
identidades fronteiriças. A possibilidade de ``não estar nem aqui, nem
ali'' permite o contato, a troca e a multiplicidade. A diáspora é uma
condição em si.

Por fim, o antropólogo e pesquisador Tim Ingold, em \emph{Estar Vivo},
nos lembra que a vida nunca é, exclusivamente, vivida neste ou naquele
lugar, ``mas sempre no caminho de um lugar ao outro''\footnote{INGOLD,
  \emph{Estar vivo: ensaios sobre movimento, conhecimento e descrição},
  p.\,218.}. A vida não é vivida dentro de um espaço, mas através dele
--- na fronteira.

Essa possibilidade de viver na fronteira, e na diáspora, faz com que os
grupos sociais habitem a terra como peregrinos, em transição, em
deslocamentos e em constantes mudanças, pois ``a experiência humana
não é situada, mas situante {[}\ldots{}{]} Ela desdobra-se não em lugares,
mas ao longo de caminhos''\footnote{\emph{Ibidem}, p.\,219.}. Assim, os
peregrinos caminham, física e mentalmente, e tecem linhas, traços, laços
e nós coletivos.

Em \emph{Estar vivo}, Tim Ingold menciona o livro \emph{Playing Dead} do
escritor canadense Rudy Wiede, que fala do povo \emph{Inuit}, habitante
da região ártica do Canadá, Alasca e Groenlândia. Ingold argumenta que,
para os \emph{Inuit}, ``assim que uma pessoa se move, ela se torna
uma linha''\footnote{\emph{Ibidem}, p.\,221.}. As linhas se transformam em
caminhos pelos quais a vida é vivida (ou escrita, como no caso dos
papéis do gueto), pois o peregrino está em contínuo movimento,
justamente por ele próprio ser o movimento. O peregrino ``é
exemplificado no mundo como uma linha de viagem {[}\ldots{}{]} É uma linha
que avança a partir da extremidade conforme ele segue
adiante''\footnote{\emph{Ibidem}.}. É por meio destes movimentos físicos, ou
destes movimentos mentais, que as linhas consolidam-se enquanto uma
possibilidade de habitar o deslocamento.

% Comentário da editora: A tradução da frase do peregrino está estranha. Ela está correta?

Ao compreender as identidades diaspóricas enquanto peregrinos que
habitam a terra, entende-se aqui que estes indivíduos ``não têm
destino final, pois onde quer que estejam, e enquanto a sua vida
perdure, há algum outro lugar aonde podem ir''\footnote{\emph{Ibidem}.}.
% Comentário da editora: é "terra" mesmo, né?

Existem duas palavras hebraicas para designar o exílio: \emph{galut} e
\emph{tefutzá}. \emph{Galut} (em português: exílio) é visto como algo
maldito, sofrido e a ser superado. \emph{Tefutzá} (em português:
dispersão), em contrapartida, é visto como um lugar em si. O exilado se
perdeu, o disperso está perdido. Ao final, o movimento de ``estar
perdido'' é uma posição política, um ato de possibilidade, de viver o
entre, o nada --- habitar a fronteira.

\begin{center}\adforn{49}\end{center}

\noindent{}Em ``Contra o espaço: lugar, movimento, conhecimento'', Tim Ingold diz
que ``espaço'' é um dos termos mais abstratos para descrever o mundo em
que vivemos, na medida em que pensar que a vida humana acontece em um
espaço --- ou seja, que ocupa uma determinada espacialidade --- está
diretamente vinculado ao que o antropólogo chama de ``lógica da
inversão''. Segundo o pesquisador, a ``lógica da inversão''
transforma as vias ao longo das quais a vida é vivida em limites nos
quais a vida é encerrada.

A diáspora é o mundo habitado, e não o mundo ocupado. A diáspora tece
lógicas a partir da construção de nós e de fios do devir, ou seja, a
partir das peregrinações e dos deslocamentos, de indivíduos, de
culturas, de tradições, de identidades, bem como de comunidades. Os
``humanos habitam a terra como peregrinos''\footnote{\emph{Ibidem}, p.
  219.} --- semeiam e cultivam nós, até que estes laços brotem para logo
se desmancharem, e serem refeitos em outras localidades. Por assim
dizer, ``lugares, então, são como nós, e os fios a partir dos
quais são atados linhas de peregrinação''\footnote{\emph{Ibidem}, p.\,220.}.

Se os movimentos diaspóricos judaicos se apropriam dos territórios do
indeterminado por meio de sua cosmogonização (os transformam em casas,
em um lugar habitável, e não espaço ocupado), e se o Deus judaico habita
o movimento, então há uma necessidade judaica de contínua e
recorrentemente se desvincular da ``abstração do
espaço''\footnote{\emph{Ibidem}, p.\,216.} para, então, habitar o lugar.

É a partir da inserção de objetos, de símbolos e signos de pertencimento
e de identificação que as comunidades judaicas criam o sentimento de
``lugar''. Há uma constante fuga do ``sentimento abstrato do
espaço''\footnote{\emph{Ibidem}.}, já que, quanto mais alta a dimensão do
espaço, mais distante o embasamento no lugar --- no ``Seu Mundo''.

Ingold argumenta que cada lugar contém, em si, lugares menores.
``Os lugares sempre se abrem para revelar outros lugares dentro de
si''\footnote{\emph{Ibidem}.}, como bonecas russas ou espaços heterotópicos.
Os lugares existem no espaço. E os lugares judaicos habitam o ``tempo'',
pois fazem do tempo o seu lugar.

O rabino e filósofo Abraham Heschel diz que ``os sábados são
nossas grandes catedrais''. O momento judaico do \emph{shabat} é um
lugar, justamente pela experiência do \emph{shabat} ser totalmente
desvinculada de um espaço. Quanto mais próximos os lugares, mais estão
atados as linhas e os laços entre a comunidade. O \emph{shabat} é
``uma volta ao lar''\footnote{\textsc{buber}, \emph{op.\,cit.}, p.\,196.}
justamente por cristalizar uma concepção de lugar.

Em diáspora, a experiência judaica vive em múltiplos e multifacetados
lugares ao mesmo tempo: ``a eternidade é conquistada por aqueles
que trocam o espaço pelo tempo''\footnote{\textsc{zevi}, \emph{op.\,cit.}, p.\,11.}.

\begin{center}\adforn{49}\end{center}

\noindent{}Os pisos das sinagogas sefaraditas\footnote{Sefaraditas
  são as comunidades judaicas provenientes de Portugal e da Espanha e
  que possuem a língua diaspórica ladino como idioma mãe.}
encontradas no Caribe (em especial no Suriname, em Curaçao e na Jamaica)
são completamente cobertos com areia proveniente de
\emph{Sefarad}\footnote{\emph{Sefarad} é o termo hebraico correspondente
  à região atual da Espanha.}. Estas sinagogas foram construídas em
meados do século \textsc{xvi}, durante o período da Inquisição Espanhola que se
alastrou por Portugal, bem como por suas colônias, e desenraizou os
judeus de seus países de origem.

Em meio à perseguição, a areia despejada sobre o piso permitia um
abafamento dos ruídos dos passos dos membros da congregação, e assim
fazia com que a vida judaica pudesse ser vivida em silêncio e em
segredo.

Para além de um movimento de resistência e de contestação, o ato de
esparramar areia sobre o piso sinagogal pode ser lido também enquanto um
lembrete de devoção à divindade e de memória ao momento de vagueio por
quarenta anos no deserto do \emph{Sinai}, antes da entrada na terra de
\emph{Canaã}. Mais que isso, a areia de \emph{Sefarad} implica uma
recordação das relações de intimidade que as comunidades
sefaraditas possuem com a sua terra natal (a Espanha) e com a
diáspora. E, afinal, o que se leva de um lugar ao outro?

O pesquisador Adolfo Roitman afirma que ``todo lugar pode ser
\emph{prima facie} uma sinagoga, implicando que o elo com Deus não reside em um
lugar fixo e determinado (como é o caso de um templo), mas que se fazia
presente onde Israel se reunia e elevava suas preces''\footnote{\textsc{roitman},
  \emph{Del Tabernáculo al Templo}, p.\,59.}. Assim sendo, a tradição
judaica entende que são os indivíduos que incorporam em si e entre si os
lugares de pertencimento ao judaísmo. Ou seja, o ``ser diaspórico'' é um
movimento compartilhado, um saber ancestral judaico e coletivo, uma
memória que acompanha as comunidades.

Estes constantes movimentos de dispersão e de transitoriedade permitem
que as comunidades judaicas e cada um dos que a compõem sejam peregrinos
e caminhem pelos territórios levando consigo seus lugares sagrados, o
``Seu Mundo'': Deus, os ritos, os símbolos e os objetos hierofânicos.
Todo novo local é demarcado e cosmogonizado.

A escritora Noemi Jaffe diz que ``ser judia é, de alguma forma,
continuar a costurar''\footnote{\textsc{jaffe}, \emph{Rebeca: A costura no
  avesso}.}. Que continuemos a tecer estes caminhos herdados e
compartilhados. Caminhos que, em algum sentido, já foram percorridos, e
serão mais uma vez, por nós. Mais uma vez apropriados, questionados,
relidos. Os caminhos judaicos estão em constante \emph{chidush}. Todos
podem e devem se apropriar, fazer deles um pouco seus.

\begin{center}\adforn{49}\end{center}

\noindent{}A intenção deste ensaio é se apresentar enquanto uma história judaica,
pois ``a vida judaica é medida pelos seus livros e permeada pela
história''\footnote{\textsc{zevi}, \emph{op.\,cit.}, p.\,7.}. Dessa forma, a tendência do
espírito judaico da diáspora é de expressar os fatos e os acontecimentos
do passado e do presente de maneiras aguçadas. As ocorrências são de tal
modo relatadas e vivenciadas que passam a dizer alguma coisa, culminando
em algo realmente dito. ``A anedota judaica afirma-se no relato de
um único incidente que aclara todo um destino. E no incidente único
exprime-se o significado da vida''\footnote{\textsc{buber}, \emph{op.\,cit.}, p.\,15.}.

E assim, na constelação destas páginas, é hora de retornar ao espelho da
página em branco e voltar a escrever palavras esparsas, signos da
diáspora --- apropriar-se da escrita, das letras e, por fim, fazer
nascer sentidos do indeterminado e de identificação.

\emph{Estos son los biervos de mis abuelas i de mis abuelos, ke ayer
cantaron mis papas, i ke hoy soi io quien los canto}\footnote{Do
  ladino: Essas são as palavras de minhas avós e de meus avôs,
  que ontem contaram aos meus pais, e que hoje sou eu quem conto.}. Hoje
sou eu quem conto, quem esboço, risco e delineio as mais diversas
linhas-palavras que querem a todo custo dizer: amanhã --- para assim
haver amanhã.

% Comentário da editora: Eu pararia aqui, e excluiria o Sheheianu.

% \noindent{}\emph{Sheheyanu v'kiyimanu Veihiguiyanu lizman hazé}\footnote{``Bendito és Tu, {[}\ldots{}{]}, que nos concedeu a vida, nos sustentou e nos permitiu chegar a este momento''.}.